\section{Conclusion}
\label{sec:conclusion}

We presented a diagnostic reproduction of \mcpilot{}~\cite{turcato2025mcpilot} in a
self-contained NumPy and PyBullet simulator, and conducted three focused studies
to characterise how the algorithm behaves outside its original calibration geometry.

Three technical contributions emerge from this work.
First, we diagnosed and corrected four non-obvious failure modes in the baseline
implementation: a broken gradient path from cost to policy parameters (T\_control
unit mismatch), cost saturation from excessive lengthscale $\ell_c$ relative to
exploration error magnitude, a wasted simulation horizon due to particles propagating
after landing, and an index-range error in the early stopping infrastructure.
Each failure was identified by instrumentation, understood through the algorithm
structure, and resolved with a minimal targeted fix.

Second, two transferable design rules emerged from Study 1.
The \emph{lengthscale-to-range scaling rule}---$\ell_s \approx 0.15 \times
\text{target\_range}$---prevents the RBF policy from becoming blind over narrow
target domains, a failure mode that is invisible from the cost curve alone.
The \emph{stratified exploration rule} partitions the speed range into equal bands
and samples one throw per band, providing seed-robust GP coverage in the same number
of exploration throws as random sampling while eliminating the max-speed local
minima that trapped random exploration on elevated configs.

Third, the release-height generalisation study demonstrates 100\% hit rates on
$z \in \{0.5, 1.0, 1.5, 2.0\}$\,m with mean landing errors below 35\,mm across all
four configurations, and 10/10 hits on Config E ($z = 0.0$\,m) after iterative
diagnosis through five progressive fixes.
These results confirm that the corrected algorithm generalises across a 2\,m range
of release heights without any per-height tuning beyond the three parameters
($\ell_M$, $T$, $\ell_s$) that are geometrically motivated.

Wind robustness is the natural next step: extending the drag model with a wind
force term (Eq.~\ref{eq:wind}) and studying whether the GP can absorb systematic
wind bias within the same five-throw exploration budget used for the static case.
Open questions include few-shot policy transfer across geometries (can a policy
trained at $z = 1.0$\,m be warm-started for $z = 2.0$\,m with fewer trials than
training from scratch?) and wind-state augmentation of the policy input, which would
transform the single-shot speed map into a full contextual policy conditioned on
measured airflow.
